% !TeX root = ../../Book5.tex
% 纲要标题：### 1.2 典型表示
% 纲要位置：工作记录/计划纲要.md，第 3729 行。
% 纲要细目（写作范围）：
% * 正则表示和置换表示
% * 一维表示及交换化
% * 对偶、张量积和外幂表示


\section{典型表示}
\label{b5:sec:0102}

抽象定义允许许多不同作用。最常用的表示往往直接来自一个有限集合：把集合元素改成向量空间的基，置换便成为线性变换。另一些表示则从已有表示通过对偶和张量构造得到。以下各项都写明作用公式，以便后面计算迹和不变子空间。


\subsection{正则表示和置换表示}
\label{b5:subsec:0102:01}

设有限群 \(G\) 作用于有限集合 \(X\)。在以 \(e_x\)、\(x\in X\)，为基的 \(k\) 空间 \(k[X]\) 上定义 \(ge_x=e_{gx}\)，得到\term[zhihuan-biaoshi]{置换表示}[permutation representation]。矩阵每行每列各有一个一，其余位置为零。\(X=G\) 且采用左平移时，得到\term[zhengze-biaoshi]{正则表示}[regular representation]，也就是 \(k[G]\) 在自身上的左乘作用。

\ProseParagraphBreak
正则表示忠实，因为 \(g\cdot1=g\)，不同基向量不相等。对一般 \(X\)，不动向量
\[
V^G=\Bigl\{v\in k[X]\mid gv=v\text{ 对所有 }g\in G\Bigr\}
\]
在同一个轨道上的系数必须相同。因此每个轨道的基向量之和给出一组不动空间的基，\(\dim_k k[X]^G\) 等于轨道数，任意特征均成立。

\begin{example}\label{b5:exm:permutation-augmentation}
设 \(S_n\) 置换 \(k^n\) 的标准基，\(n\ge2\)。令
\[
L=k(e_1+\cdots+e_n),\qquad
W=\Bigl\{\sum_i a_ie_i\mid\sum_i a_i=0\Bigr\}.
\]
两者都是子表示。当 \(n\ne0\) 于 \(k\) 时，\(k^n=L\oplus W\)；当 \(\operatorname{char}k\mid n\) 时，却有 \(L\subseteq W\)。
\end{example}
\begin{proof}
置换不改变系数和，故两者不变。常数向量 \(a\sum_i e_i\) 的系数和为 \(na\)。若 \(n\) 可逆，交为零且维数相加为 \(n\)，得到直和；否则整个 \(L\) 都在 \(W\) 中。
\end{proof}

\subsection{一维表示及交换化}
\label{b5:subsec:0102:02}

一维 \(k\) 表示就是群同态 \(\lambda:G\to k^\times\)，称为\term[yiwei-biaoshi]{一维表示}[one-dimensional representation]。像群交换，所以 \(\lambda(ghg^{-1}h^{-1})=1\)。记 \(G'=[G,G]\) 为所有交换子生成的子群，则 \(\lambda\) 唯一经由 \(G/G'\) 分解。反过来，\(G/G'\to k^\times\) 的同态与商映射复合就给出一维表示。这直接说明一维表示只能观察群的交换化。

\ProseParagraphBreak
在复数域上，\(C_m=\langle r\rangle\) 的一维表示为
\[
\lambda_j(r)=\exp(2\pi\mathrm i j/m),\qquad 0\le j<m.
\]
在实数域上，有限群的一维表示的值只能为 \(\pm1\)；阶为奇数的群因而只有平凡实一维表示。同一个群换底域以后，可用的一维表示会改变。

\ProseParagraphBreak
若 \(\lambda\) 为一维表示，\(\rho^\lambda(g)=\lambda(g)\rho(g)\) 仍满足表示乘法法则，称为用 \(\lambda\) 扭曲。子空间在扭曲前后不变的条件相同，因此不可约性不变。例如对称群的符号表示可将一个不可约表示变成另一个，而不改变次数。

\subsection{对偶、张量积和外幂表示}
\label{b5:subsec:0102:03}

设 \(V,W\) 是群 \(G\) 在域 \(k\) 上的有限维表示。对偶空间 \(V^*=\operatorname{Hom}_k(V,k)\) 上定义
\begin{equation}\label{b5:eq:dual-tensor-action}
(g\varphi)(v)=\varphi(g^{-1}v),\qquad
g(v\otimes w)=gv\otimes gw.
\end{equation}
前式给出\term[duiou-biaoshi]{对偶表示}[dual representation]，后式给出\term[zhangliang-biaoshi]{张量积表示}[tensor product representation]。逆元不可省略：计算
\[
g(h\varphi)(v)=\varphi(h^{-1}g^{-1}v)=\Bigl((gh)\varphi\Bigr)(v)
\]
才得到左表示的次序。对偶矩阵是 \(\rho(g^{-1})^{\mathsf T}\)，没有共轭；复共轭只会在另行选择 Hermitian 内积后出现。

\ProseParagraphBreak
外幂上的作用由
\[
g(v_1\wedge\cdots\wedge v_r)=gv_1\wedge\cdots\wedge gv_r
\]
定义，因为线性变换保持交替关系，故能从张量幂下降到商空间。特别地，最高外幂是一维表示 \(g\mapsto\det\rho(g)\)。在特征二，外代数仍按 \(v\wedge v=0\) 定义，不能仅用反交换公式替代交替条件。

\begin{proposition}\label{b5:prop:hom-tensor-invariants}
设 \(k\) 为域，\(G\) 为群，\(V,W\) 为 \(G\) 的有限维 \(k\) 表示。在 \(\operatorname{Hom}_k(V,W)\) 上置 \((gT)(v)=gT(g^{-1}v)\)，其中 \(g\in G\)、\(T\in\operatorname{Hom}_k(V,W)\)、\(v\in V\)，则自然映射
\[
W\otimes_kV^*\longrightarrow\operatorname{Hom}_k(V,W),
\qquad w\otimes\varphi\longmapsto\Bigl(v\mapsto\varphi(v)w\Bigr)
\]
是表示同构，且其不动空间恰为 \(\operatorname{Hom}_G(V,W)\)。
\end{proposition}
\begin{proof}
选 \(V\) 的基 \(v_i\) 及对偶基 \(\varphi_i\)，逆映射为 \(T\mapsto\sum_iT(v_i)\otimes\varphi_i\)，与所选基无关，因为它是给定自然映射的唯一逆。两侧作用于纯张量后均将 \(v\) 送到 \(\varphi(g^{-1}v)gw\)，故同构与作用相容。最后 \(gT=T\) 等价于 \(T(gv)=gT(v)\)，即交织条件。
\end{proof}
